\chapter{Kirchengehen}
\quellseite{9}

Zwei Bekannte gingen mit einander in die Kirche. „Was nützt es eigentlich“, sagte der Jüngere, „so oft in die Kirche zu gehen?“

„Was nützt es“, erwiderte der andere, „so oft zu essen?“

„O das ist doch etwas ganz anderes; durch die Speisen erhalte ich mein Leben und meine Kraft“, war die Antwort.

„Die beiden Dinge sind nicht so verschieden, wie du denkst“, sagte nun wieder der Aeltere. „Was die Speise für den Leib ist, das ist das Wort Gottes für die Seele.“

„Aber“, wendete jener ein, „woher kommt es denn, daß so viele nach dem Worte Gottes thatsächlich kein Verlangen tragen, während doch jedermann essen möchte?“

„Du täuschest dich“, war die Antwort, „nicht jeder will essen. Geh nur in ein Lazarett und sieh die Kranken an! Allerdings ist es das Naturgemäße, daß ein Mensch Appetit hat; hat er keinen, so ist er krank. Und wenn so viele nach dem Worte Gottes kein Verlangen tragen, so ist das eben ein sehr ernstes Zeichen weitverbreiteter und schwerer innerer Erkrankung.“

In der That, wenn es mit unserer Seele richtig steht, dann kann sie die Erbauung aus Gottes Wort nicht entbehren. Ist sie aber verhärtet, so fehlt ihr der Hunger nach geistlicher Nahrung, ja, sie geht ihr aus dem Wege. Aber weit entfernt, daß dies ein Zustand der Gesundheit und Kraft wäre, ist es vielmehr ein kranker und gefährlicher Zustand.
